% 题干转写：完成
% 解析转写：完成
% 图形核对：完成
\section{第40届复赛第4题：编木拱桥的几何与静力平衡}
\subsection{题干}
北宋《清明上河图》中的编木拱桥以及类似实物模型，都利用杆件间的支持力与静摩擦力保持结构稳定。现取由三个单元组成的左右对称简化模型，只研究竖直平面内的平衡。每对长杆等效为质量为 $M$ 的杆，结构由短杆 $A,A',B$ 和长杆 $C,D,D'$ 组成。各短杆质量为 $m$、半径为 $R$；各长杆长度为 $L$；短杆 $B$ 表面光滑，其余杆间静摩擦因数为 $\mu$，杆与地面间的静摩擦因数为 $\mu'$。$D,D'$ 与水平面夹角均为 $\theta$，如图\ref{fig:cpho40q04}。
\begin{figure}[htbp]\centering
\begin{tikzpicture}[scale=0.62]
  \begin{scope}
    \draw[thick] (-2.8,0) -- (2.8,0);
    \foreach \x in {-2.6,-2.2,...,2.6} {\draw[gray] (\x,0) -- +(0.24,-0.16);}
    \draw[very thick] (-2.2,0) .. controls (-1.0,2.7) and (1.0,2.7) .. (2.2,0);
    \draw[thick] (-2.2,0) -- (-0.45,1.9) -- (1.25,0);
    \draw[thick] (-1.25,0) -- (0.45,1.9) -- (2.2,0);
    \node[below] at (0,-0.6) {图4a：编木拱桥结构示意};
  \end{scope}
  \begin{scope}[xshift=7cm]
    \draw[thick] (-2.8,0) -- (2.8,0);
    \foreach \x in {-2.6,-2.2,...,2.6} {\draw[gray] (\x,0) -- +(0.24,-0.16);}
    \draw[very thick] (-2.3,0) -- (-0.5,2.1) -- (1.3,0);
    \draw[very thick] (-1.3,0) -- (0.5,2.1) -- (2.3,0);
    \draw[very thick] (-2.0,0.55) -- (1.95,0.55);
    \draw[very thick] (-1.4,1.15) -- (1.4,1.15);
    \node[below] at (0,-0.6) {图4b：多单元模型结构示意};
  \end{scope}
\end{tikzpicture}

\caption{依据题图4a、4b转绘的编木拱桥与多单元模型结构示意。}\label{fig:cpho40q04context}
\end{figure}
\begin{figure}[htbp]\centering
\begin{tikzpicture}[>=Stealth,scale=0.75]
  \draw[thick] (-4,0) -- (4,0);
  \draw[very thick] (-3.6,0.1) -- (-0.2,3.5) -- (1.2,2.1) -- (3.6,0.1);
  \draw[very thick] (-2.5,0.1) -- (0.6,3.2) -- (2.2,1.6);
  \draw[very thick] (-3.1,1.35) -- (3.1,1.35);
  \draw[very thick] (-2.45,0.32) -- (-0.2,2.57);
  \draw[dashed] (-2.62,0.15) -- (-0.37,2.40);
  \node[left] at (-3.65,0.15) {$D$}; \node[right] at (3.65,0.15) {$D'$};
  \node[left] at (-3.15,1.35) {$C$}; \node at (-1.45,1.06) {$A$};
  \node at (2.65,1.06) {$A'$}; \node at (-0.15,1.72) {$B$};
  \draw[<->] (-2.78,0.25) -- (-1.87,1.16) node[midway,left] {$a$};
  \draw[<->] (-2.49,0.54) -- (-0.67,2.36) node[midway,right] {$b$};
  \draw[->] (-3.15,0.25) arc[start angle=35,end angle=0,radius=0.55];
  \node at (-2.62,0.37) {$\theta$};
\end{tikzpicture}

\caption{第40届复赛第4题图4c的TikZ重绘：三单元编木拱桥侧视结构。}\label{fig:cpho40q04}
\end{figure}
\begin{enumerate}
\item 若 $D$ 杆下端至与 $A$ 杆接触点的长度为 $a$，求其至与 $B$ 杆接触点的长度 $b$。
\item 分析 $A,B,C,D$ 四杆受力，画出 $C,D$ 的受力图，并列出平衡方程。
\item 当 $M=6m$、$\theta=45^\circ$、$L=30\,\mathrm{cm}$、$a=16\,\mathrm{cm}$、$R=2(\sqrt2-1)\,\mathrm{cm}$ 时，求保持平衡的 $\mu,\mu'$ 条件。
\end{enumerate}
\subsection{参考解答整理}
\paragraph{（1）接触位置的几何关系。}圆柱杆接触的几何关系给出
\[b=a+4R\cot\frac\theta2.\]
\paragraph{（2）受力与平衡方程。} $C,D$ 两杆的受力关系以图\ref{fig:cpho40q04forces}表示。对 $A,B,C,D$ 分别列平衡条件，可写为
\[
N_1+mg=N_2\cos\theta+f_2\sin\theta,\quad
f_1+f_2\cos\theta=N_2\sin\theta,\quad f_1=f_2,
\]
\[N_3=2N_4\cos\theta+mg,\qquad 2N_1=N_3+Mg,
\]
\[
N_5+N_4\cos\theta=N_2\cos\theta+f_2\sin\theta+Mg,
\quad f_5+N_2\sin\theta=N_4\sin\theta+f_2\cos\theta.
\]
以 $D$ 杆与地面的接触点取力矩，
\[
N_2a+Mg\left(\frac L2\cos\theta-R\sin\theta\right)=N_4b+2Rf_2.
\]
\begin{figure}[htbp]\centering
\begin{tikzpicture}[>=Stealth,scale=0.72]
  \begin{scope}
    \draw[very thick] (-1.65,0) -- (1.65,0);
    \draw[thick] (-0.95,-0.36) circle (0.28);
    \draw[thick] (0.95,-0.36) circle (0.28);
    \draw[thick] (0,0.36) circle (0.28);
    \draw[->,thick] (-0.95,0) -- (-0.95,0.92) node[above] {$N_1$};
    \draw[->,thick] (0.95,0) -- (0.95,0.92) node[above] {$N_1$};
    \draw[->,thick] (-1.25,0) -- (-1.83,0) node[left] {$f_1$};
    \draw[->,thick] (1.25,0) -- (1.83,0) node[right] {$f_1$};
    \draw[->,thick] (0,0) -- (0,-0.88) node[right] {$N_3,\ Mg$};
    \node[below] at (0,-1.12) {$C$ 杆受力示意};
  \end{scope}
  \begin{scope}[xshift=4.3cm]
    \draw[thick] (-1.45,0) -- (1.45,0);
    \draw[very thick] (-0.82,0.08) -- (1.06,1.96);
    \draw[thick] (-0.35,0.56) circle (0.28);
    \draw[thick] (0.62,1.53) circle (0.28);
    \draw[->,thick] (-0.82,0.08) -- (-0.82,0.96) node[left] {$N_5$};
    \draw[->,thick] (-0.82,0.08) -- (-0.06,0.08) node[below] {$f_5$};
    \draw[->,thick] (-0.35,0.56) -- (0.10,0.11) node[right] {$N_2$};
    \draw[->,thick] (-0.82,0.08) -- (-1.28,-0.38) node[left] {$f_2$};
    \draw[->,thick] (0.62,1.53) -- (0.13,2.02) node[left] {$N_4$};
    \draw[->,thick] (0.10,1.00) -- (0.10,0.27) node[right] {$Mg$};
    \node[below] at (0,-0.62) {$D$ 杆受力示意};
  \end{scope}
\end{tikzpicture}

\caption{依据题解图4a、4b重绘的 $C$ 杆与 $D$ 杆受力示意。}\label{fig:cpho40q04forces}
\end{figure}
\paragraph{（3）代入数值。}先得 $b=24\,\mathrm{cm}$。由上述方程组，
\[N_1=10.427mg,\qquad f_1=f_2=4.733mg,\qquad N_5=\frac{21}{2}mg,\qquad f_5=2.194mg.\]
杆件接触处需满足 $f_1\leq\mu N_1$、$f_2\leq\mu N_2$，地面处需满足 $f_5\leq\mu'N_5$。因此保持平衡的条件为
\[\mu\geq0.454,\qquad \mu'\geq0.209.\]
\subsection{转写依据}
题干和解析分别依据 \texttt{CPHO40-Q04\_transcription\_v1.md} 与 \texttt{CPHO40-Q04\_formula\_check\_v1.csv}；图4a至4c及题解受力图的结构和力标注均据相应图件核对记录重绘。图4a、4b原为历史/实物照片，图册保留其结构信息并以TikZ示意替代照片，不嵌入原始图片。
